\documentclass[12pt]{article}
\usepackage[utf8]{inputenc}
\usepackage[T1]{fontenc}
\usepackage[margin=1in]{geometry}
\usepackage{setspace}
\usepackage{xcolor}
\usepackage[colorlinks=true, urlcolor=blue!60!black, linkcolor=blue!60!black]{hyperref}
\setlength{\parindent}{2em}
\setlength{\parskip}{1.25em}
\renewcommand{\baselinestretch}{1.0}
\begin{document}
\begin{center}
{\LARGE \textit{Luther v. Borden}, 48 U.S. (7 How.) 1 (1849)\par}
\vspace{0.8em}
{\large Case Brief --- Four-Page Edition\par}
\end{center}
\vspace{1em}

\section{Memory Jogger}

Federal courts lack authority to decide which of two rival state governments is legitimate under the Guarantee Clause, because recognizing a state's lawful government is a nonjusticiable political question committed to Congress and the President (\href{https://supreme.justia.com/cases/federal/us/48/1/}{48 U.S.\ 1}).

\section{Detailed Case Facts}

Rhode Island still operated under King Charles II's 1663 colonial charter, which had no amendment mechanism and let the General Assembly restrict voting to freeholders, excluding a growing majority of adult men. After petitions for suffrage reform failed, reformers convened an unauthorized ``People's Convention'' in 1841, drafted the ``People's Constitution'' extending suffrage to adult white male citizens, held an extralegal ratification vote, and in April 1842 elected Thomas Wilson Dorr governor along with a rival legislature.

The charter government under Governor Samuel Ward King refused to yield. After Dorr tried to seize the Providence arsenal and gathered armed supporters at Chepachet, the charter Assembly declared martial law on June 25, 1842, and King called out the militia and asked President Tyler for help. Tyler sent no troops but recognized the charter government and signaled readiness to act, and Dorr's forces dispersed. On June 29, 1842, militiaman Luther M.\ Borden and others, under orders to arrest insurgents, broke into Martin Luther's home in Warren to arrest him for aiding Dorr. A new constitution adopted through the charter government took effect in May 1843.

\section{Procedural History}

Luther, having moved to Massachusetts, sued Borden and the other militiamen for trespass in the federal circuit court. The defendants justified their entry under martial law. At trial, Justice Story, acting \textit{pro forma} by agreement, excluded Luther's evidence that the People's Constitution had lawfully displaced the charter government and instructed the jury that the charter government's laws justified the defendants. The jury found for the defendants, and Luther brought a writ of error. A companion suit by his mother came up on a certificate of division. The Court decided both on January 3, 1849.

\section{Judicial Votes}

Affirmed, \textbf{5--1}.
\begin{itemize}
\item \textbf{Majority:} Chief Justice Taney, joined by Justices McLean, Wayne, Nelson, and Grier.
\item \textbf{Dissenting in part:} Justice Woodbury, agreeing that legitimacy was a political question but dissenting on martial law.
\item \textbf{Not participating:} Justices Catron, Daniel, and McKinley.
\end{itemize}

\section{Holding}

Federal courts cannot adjudicate which rival regime is a state's lawful government or whether a state government satisfies the Guarantee Clause of Article IV, Section 4. That determination belongs to Congress, in seating state delegations, and to the President, in deciding whether to intervene against domestic violence under the Act of 1795, and their recognition binds the courts. A state may also declare martial law to suppress an insurrection too strong for civil authority, and officers may arrest suspected insurgents and enter homes on reasonable grounds using no more force than necessary. Taney's suggestion that a permanent military government would violate the Guarantee Clause is dicta, and Tara Leigh Grove notes that his sweeping nonjusticiability statements went beyond the pleadings. The companion certificate was dismissed.

\section{Analysis of Opinions}

\textbf{Majority (Taney).} Taney rested on four pillars. \emph{Practical chaos:} voiding the charter government retroactively would nullify every law, tax, salary, and judgment of 1842--1843 and expose officers to liability. \emph{State-court deference:} Rhode Island's courts, including in Dorr's treason trial, had upheld the charter government, and a state court could not declare its own creator illegitimate. \emph{No manageable standards:} a court cannot set voter qualifications or count thousands of voters, and different juries could reach conflicting verdicts on ratification. \emph{Textual commitment:} the Guarantee Clause entrusts the guarantee to the political branches; Congress acts by seating representatives, and under \textit{Martin v.\ Mott} the President's recognition under the 1795 Act was conclusive. On the merits, martial law was a legitimate measure of self-defense, though officers remain answerable for oppressive or willful excess.

\textbf{Woodbury (dissenting in part).} Woodbury agreed that courts operate beneath constitutions, not above the people in forming them, warning that unelected judges with such power would form an irresponsible oligarchy. But martial law over civilians was a judicial question involving private rights. Distinguishing military law, which governs soldiers, from martial law, which suspends civil law for everyone, and invoking Magna Carta and the Petition of Right, he argued that no legislature could subject civilians to warrantless military arrest while courts were open; officials acting in emergencies should seek indemnity afterward rather than advance sanction for ``the reign of the bayonet.''

\section{Examples: Future Applications}

\begin{enumerate}
\item \textbf{Same-Side: Direct Democracy.} A challenge to lawmaking by ballot initiative as un-republican is a political question, as in \textit{Pacific States Telephone \& Telegraph Co.\ v.\ Oregon} (1912).
\item \textbf{Same-Side: Partisan Gerrymandering.} Echoing \textit{Luther}'s concern about manageable standards, \textit{Rucho v.\ Common Cause} (2019) held partisan gerrymandering claims nonjusticiable.
\item \textbf{Opposite-Side: Malapportionment.} \textit{Baker v.\ Carr} and \textit{Reynolds v.\ Sims} distinguished \textit{Luther}: Equal Protection claims are justiciable because ``one person, one vote'' is manageable and does not question a government's legitimacy.
\item \textbf{Opposite-Side: Excessive Force or Military Trials.} Militia officers who destroy a peaceful journalist's press out of animus and try her by military tribunal while courts are open remain answerable, even under Taney's opinion.
\item \textbf{Fence-Sitter: Commandeering.} In \textit{New York v.\ United States} (1992), Justice O'Connor suggested ``perhaps not all claims under the Guarantee Clause present nonjusticiable political questions'' but did not decide the point.
\end{enumerate}

\section{Critique}

\textbf{Originalist and textualist views.} Ryan C.\ Williams and others note that the Clause assigns the duty to ``[t]he United States,'' not only Congress, leaving room for judicial enforcement. Tara Leigh Grove argues that early political-question cases involved deference to specific political-branch determinations, not a jurisdictional bar. Other formalists defend \textit{Luther} as restraining unelected judges.

\textbf{Progressive and doctrinal views.} Erwin Chemerinsky argues that \textit{Luther} erased the Guarantee Clause, abdicating the \textit{Marbury} duty and leaving democratic deficits without remedy. Historians note that nonjusticiability shielded slave states from scrutiny. Woodbury's martial-law dissent is widely praised as a prescient defense of civil liberties.

\section{Key Quotations}

\begin{quote}
``Under this article of the Constitution, it rests with Congress to decide what government is the established one in a State. For as the United States guarantee to each State a republican government, Congress must necessarily decide what government is established in the State before it can determine whether it is republican or not.''
\end{quote}
\noindent --- Chief Justice Taney, 48 U.S.\ (7 How.)\ at 42.

\begin{quote}
``If the judicial power extends so far, the guarantee contained in the Constitution of the United States is a guarantee of anarchy, and not of order.''
\end{quote}
\noindent --- Chief Justice Taney, 48 U.S.\ (7 How.)\ at 43.

\begin{quote}
``And unquestionably a State may use its military power to put down an armed insurrection too strong to be controlled by the civil authority. The power is essential to the existence of every government, essential to the preservation of order and free institutions, and is as necessary to the States of this Union as to any other government.''
\end{quote}
\noindent --- Chief Justice Taney, 48 U.S.\ (7 How.)\ at 45.

\end{document}
